\section{Metrics}
\label{sec:metrics}

\subsection{Accuracy computation}
We compute F\textsubscript{1} from true positives (TP), false positives (FP), and false negatives (FN) defined at the row–value level. 
For Multi rows, each predicted value is treated independently (“unrolled”) to avoid crediting extra unattested values. 
LLM responses that cannot be mapped to any schema row or valid value type are marked invalid and excluded from scoring. 
When a predicted value is type‑coercible (\textit{e.g.}, the string ``4'' for a Numeric row), we normalize it before comparison. 
All F\textsubscript{1} results are averaged over three independent runs, with coefficients of variation computed across runs.

\subsection{Latency and token accounting}
We measure per‑transcript latency and token usage separately for (i) prompt/input tokens and (ii) completion/output tokens.
For the baseline (no segmentation), latency is simply the extraction call duration.
For SchemaRAG, we measure (a) the segmentation call and (b) the $n$ extraction calls; because extraction calls run in parallel, the effective extraction latency is the maximum of these $n$ durations.
Final reported latency and token counts are averaged over three runs for each experimental condition, also with coefficients of variation.

\subsection{Significance testing}
To assess whether differences in F\textsubscript{1}, latency, or token usage are statistically reliable, we apply the Wilcoxon signed-rank test \cite{wilcoxon1992individual} to per-transcript metrics, averaged over three runs. 
This non-parametric test is chosen because Shapiro–Wilk tests \cite{shapiro1965analysis} indicate that per-transcript distributions significantly deviate from normality, often exhibiting skewness or heavy tails.
Unless otherwise stated, we treat $p < 0.05$ as statistically significant.
